% Generated draft: requires mathematical and source review.
\documentclass[11pt]{article}
\usepackage[margin=25mm]{geometry}
\usepackage{fontspec}
\setmainfont{Latin Modern Roman}
\usepackage{amsmath,amssymb,mathtools,unicode-math}
\setmathfont{Latin Modern Math}
\usepackage{enumitem,needspace,xurl,hyperref}
\setlength{\parindent}{0pt}
\setlength{\parskip}{5pt}
\setlength{\emergencystretch}{3em}
\newcommand{\N}{\mathbb{N}}
\newcommand{\Z}{\mathbb{Z}}
\newcommand{\Q}{\mathbb{Q}}
\newcommand{\R}{\mathbb{R}}
\newcommand{\C}{\mathbb{C}}
\newcommand{\F}{\mathbb{F}}
\newcommand{\A}{\mathbb{A}}
\newcommand{\PP}{\mathbb{P}}
\newcommand{\E}{\mathbb{E}}
\newcommand{\eps}{\varepsilon}
\newcommand{\dd}{\,\mathrm{d}}
\begin{document}
\section*{UnsolvedMath ID: 20002288}
\textbf{Source problem code:} catalogue entry
\subsection{Definitions and mathematical statement}
Pheidas's Question \(4\) asks whether the constant subfield \(\Q\subset\Q(T)\) is existentially definable in the ring language with the named transcendental element \(T\). Explicitly, does there exist an existential formula \(\delta(X)\) in that language such that
\[
 \Q(T)\models\delta(a)\quad\Longleftrightarrow\quad a\in\Q
 \qquad(a\in\Q(T))?
\]
The source says ordinary first-order definability is known even without naming \(T\), and elliptic curves give existentially definable dense subsets of \(\Q\). It gives no topology for dense and no assertion that density alone defines the full field.

The aid reports the unconditional question still open in primary and author-maintained literature checked through 10 August 2026. It attributes a 2025 positive conditional result to Garcia--Fritz and Pasten, assuming two elliptic-surface conjectures. Their exact statements and full citation are absent; no unconditional definition follows from this report.

The aid's generic-section obstruction concerns a fixed polynomial-witness presentation
\[
 D=\{a\in\Q(T):\exists y_1,\ldots,y_m\in\Q(T),\
       F_i(T,a,y_1,\ldots,y_m)=0\ (1\leq i\leq r)\},
\]
with \(F_i\in\Q[T,X,Y_1,\ldots,Y_m]\), after clearing coefficient denominators. Suppose there is a rational witness section
\[
 y_j=R_j(T,X)\in\Q(T,X)
\]
satisfying all equations identically at the generic \(X\). Then, for all but finitely many \(c\in\Q\), substitution \(X=T+c\) makes every denominator a nonzero rational function of \(T\), and the specialized witnesses lie in \(\Q(T)\). Hence
\[
 T+c\in D\qquad\text{for all but finitely many }c\in\Q.
\]
Since these elements are nonconstant, such a presentation cannot define exactly \(\Q\). The finite exceptional set concerns identically vanishing specialized denominators, not zeros at individual complex values of \(T\). The obstruction does not say that a formula defining all rational constants must possess a rational generic section.

A separate amplification observation is: if an existentially definable set \(D\subseteq\Q\) is a finite additive basis, meaning that for some fixed \(N\geq1\),
\[
 \Q=\underbrace{D+\cdots+D}_{N\text{ terms}},
\]
then all of \(\Q\) is existentially definable in \(\Q(T)\) by taking \(N\) witnesses in \(D\) and their sum. Variants using a stated finite signed sum have the same form. The finite-basis property is extra arithmetic information; it does not follow merely from the source's density assertion. Neither this amplification nor the section obstruction settles the unconditional problem.
\subsection{Short English statement}
The source asks for an existential definition of \(\Q\) inside \(\Q(T)\). The aid reports a conditional 2025 result and an unconditional open status through August 2026. A rational generic witness section cannot give the desired definition, because it would also admit almost every nonconstant \(T+c\); density of a definable constant subset is not by itself enough.
\subsection{Sources}
\begin{itemize}
\item Local UnsolvedMath registry, permanent ID \(20002288\), source code \texttt{catalogue entry}; canonical collection \texttt{aim-logic-notes.json}, notes[63].
\item AIM workshop ``Definability and decidability problems in number theory,'' Question \(4\), Pheidas: \url{https://aimath.org/pastworkshops/definabilityinntproblems.pdf}. The trailing extracted ``1'' supplies no mathematical condition. No PDF was opened in this session.
\item Garcia--Fritz and Pasten (2025), conditional positive result assuming two elliptic-surface conjectures, as reported by the aid. No full title, identifier, exact conjectures or authoritative publication data are supplied. The unconditional status qualification is checked only through 10 August 2026.
\item Supplied research attempt \(1\), \texttt{partial\_result}, run status ``open,'' rigor label ``proved,'' novelty confidence ``low.'' Registry status \texttt{partially\_solved}, classification \texttt{PARTIAL-PROGRESS}, updated 14 August 2026. Original remarks and literature are empty; no proposer-field value or year is supplied.
\end{itemize}
\end{document}
